% Source: https://artofproblemsolving.com/community/c3h1927144_math_problems
% Stable identity: aops-topic-1927144-q03
Let the ant start at $P$ and go to the opposite vertex $Q$. No single face contains both $P$ and $Q$, so any path on the surface crosses from a face containing $P$ to a face containing $Q$, passing through a common edge of the two faces. Unfold those two adjacent faces into a $1\times2$ rectangle; the path becomes a path in the rectangle from one corner to the opposite corner, which is at least as long as the straight segment.

\begin{asy}
usepackage("amsmath");
size(3.2cm);
defaultpen(fontsize(10pt));
draw((0,0)--(1,0)--(1,2)--(0,2)--cycle,linewidth(1));draw((0,1)--(1,1),linewidth(1));
draw((0,0)--(1,2),red+linewidth(1));
dot((0,0));dot((1,2));
label("$P$",(0,0),SW);label("$Q$",(1,2),NE);
label("$1$",(0.5,0),S);label("$1$",(1,0.5),E);label("$1$",(1,1.5),E);
label("front",(0.72,0.35),fontsize(7pt));label("top",(0.3,1.72),fontsize(7pt));
label("$\sqrt5$",(0.24,1.33),red);
label("\mbox{}",(min(currentpicture,user=true)+max(currentpicture,user=true))/2,black); // keeps the text after the figure black when LaTeX typesets the labels inline
\end{asy}

That segment has length $\sqrt{1^2+2^2}=\sqrt5$, and it is achieved: the straight line in the net crosses the shared edge at its midpoint and folds back to a path on the cube. (A path that wanders through more faces can also be unfolded along the faces it visits, and the same estimate holds for the part between the last face touching $P$ and the first face touching $Q$.) The shortest distance is $\boxed{\sqrt5}$.
